\section*{Chapter \ref{ch:nw:cross-region-crypto-networks} --- Cross-Region Crypto Networks}

\begin{solution}{exo:nw:cross-region-crypto-networks:1}
Against a fibre round-trip floor of \qty{91.3}{\milli\second}: $224.3/91.3 = 2.46$ for the backbone and $136.5/91.3 = 1.50$ for the specialist.
\end{solution}

\begin{solution}{exo:nw:cross-region-crypto-networks:2}
BGP chooses among routes by the policies of the networks that announce them (business relationships, preferences, path length in \omterm{def:nw:where-the-crypto-venues-live:as}{autonomous systems}),
not by measured delay; it does not know the latency of a path at all.
\end{solution}

\begin{solution}{exo:nw:cross-region-crypto-networks:3}
Frankfurt $224.3 - 136.5 = 87.8$~ms; London $210.0 - 139.6 = 70.4$~ms; Stockholm $245.3 - 124.3 = 121.0$~ms of round trip.
\end{solution}

\begin{solution}{exo:nw:cross-region-crypto-networks:4}
It correlates their delays and their failures: congestion or a fault on the shared cable hits both, so the minimum of the two is often no better
than either, and racing saves little in exactly the events it was meant to cover.
\end{solution}

\begin{solution}{exo:nw:cross-region-crypto-networks:5}
Safe: messages whose repetition has no effect (cancels, which fail harmlessly once the order is gone; queries; heartbeats), and orders that rest
long enough that every copy arrives while the first is open. Unsafe: orders that can fill on arrival (taking orders, immediate-or-cancel), whose late
copies are accepted as new orders once the first has filled.
\end{solution}

\begin{solution}{exo:nw:cross-region-crypto-networks:6}
Next to the venue whose prices it must react to fastest, usually the one it quotes or takes on, with the other market's data carried to it on the
fastest route; if it hedges on both, an engine in each region with a fast link between them. The \qty{100}{\milli\second}-scale round trip means the
engine must be where the race is.
\end{solution}

\begin{solution}{exo:nw:cross-region-crypto-networks:7}
Between $\rho = 0.46$ (\qty{110}{\micro\second}) and $\rho = 0.48$ (\qty{98}{\micro\second}): above about 0.48 the third route saves less than
\qty{100}{\micro\second} at the 99th percentile.
\end{solution}

\begin{solution}{exo:nw:cross-region-crypto-networks:8}
A venue that rejects duplicates only while the first order is open accepts a late copy once the first has filled: for taking orders, racing on three
routes can mean three executions. It is free latency only for messages the venue treats idempotently.
\end{solution}

\begin{solution}{pb:nw:cross-region-crypto-networks:1}
\textbf{Part I.}
\begin{enumerate}
\item \qty{5311}{\kilo\metre}; round-trip floors of \qty{35.4}{\milli\second} in vacuum and \qty{51.8}{\milli\second} in fibre.
\item 1.26 for the specialist, 1.31 for the backbone.
\item $65.4 - 51.8 = \qty{13.6}{\milli\second}$ of round trip.
\item Because nobody published it: it is the chapter's assumption for a path chosen by BGP, slower and with a heavier tail.
\end{enumerate}
\textbf{Part II.}
\begin{enumerate}[resume]
\item 1\,623, 771 and \qty{53}{\micro\second}.
\item 2\,141, 861 and \qty{53}{\micro\second}.
\item At 10~USD per microsecond: 5\,182, 902 and 0~USD a month.
\item It lowers it too, by \qty{536}{\micro\second} at $\rho = 0$ with two routes and 290 at 0.5, because the backbone sometimes beats the
  specialist; but the gain is largest in the tail.
\end{enumerate}
\textbf{Part III.}
\begin{enumerate}[resume]
\item Two: every late copy.
\item 1.57, 0.50 and 0.07 on average.
\item Cancels, amendments where the venue treats them idempotently, and queries; not taking orders.
\item That a client order identifier be rejected for the life of the session (or a day), not only while the order is open.
\end{enumerate}
\textbf{Part IV.}
\begin{enumerate}[resume]
\item \emph{Named result}: racing two routes saves \qty{1623}{\micro\second} of 99th percentile when they are independent, 771 at $\rho = 0.5$ and 53 at
  0.9; a third saves 518, 90 and 0 more, so at 10~USD per microsecond a month racing it stops paying above 5\,182, 902 and 0~USD a month.
\item The correlation $\rho$.
\item From logged copies: send probes (or every message) on all routes with timestamps and compute the correlation of the routes' delays, with care for
  the tails.
\item It jumps, if both routes share the cut cable's replacement path, or falls, if one route moves to a separate path; re-measure after every event.
\item Racing the backbone against a specialist \qty{88}{\milli\second} faster saves nothing unless the specialist fails; use the backbone as a fall-back,
  not a racer.
\item Everything is added to it: a slow instance at either end wastes the fastest route, so the lottery comes first.
\item In the cross-region rows: routes, prices, the racing rule per message type, and the venues' duplicate rules.
\item When the routes fail independently, the tail matters to the strategy, and the venue will not execute a late copy twice.
\end{enumerate}
\end{solution}

\begin{solution}{iq:nw:cross-region-crypto-networks:1}
The backbone's path is chosen for capacity and resilience, not latency, and may go a long way round; add the equipment and the routing between the
region's zones and the backbone. The published figures put it at 2.5 times the fibre floor, where a specialist runs at 1.5.
\iqlookfor{\omterm{def:nw:cross-region-crypto-networks:infl}{Path inflation}; routing by policy; the floor as reference.}
\end{solution}

\begin{solution}{iq:nw:cross-region-crypto-networks:2}
When the strategy's edge depends on seeing another region's prices sooner than competitors (cross-venue arbitrage, hedging), when the saving is tens
of milliseconds as between Tokyo and Europe, and when the traded volume pays for it; not when the backbone is already close to the floor, as to
Singapore.
\iqlookfor{Value of latency against cost; where the backbone is bad; measurement.}
\end{solution}

\begin{solution}{iq:nw:cross-region-crypto-networks:3}
The minimum of the two delays: a lower median and a much lower tail if the routes fail independently, little if they are correlated. It depends on the
correlation, on each route's tail, and on whether the receiver can discard the second copy safely.
\iqlookfor{Order statistics; correlation; idempotency.}
\end{solution}

\begin{solution}{iq:nw:cross-region-crypto-networks:4}
An identifier attached to every copy of one request so that the receiver acts once. It fails if the receiver forgets the identifier too soon (for
example, once the order has filled), if the sender reuses identifiers, or if different copies carry different keys.
\iqlookfor{Receiver-side memory; lifetime of keys; sender discipline.}
\end{solution}

\begin{solution}{iq:nw:cross-region-crypto-networks:5}
It makes cross-venue strategies local: an engine must sit next to the venue it trades, see the other market a tenth of a second late, and trade only
signals that survive that delay; fast arbitrage between regions is a contest of routes, won by those on the best private line.
\iqlookfor{Engine placement; signal horizon; the value of private lines.}
\end{solution}

\begin{solution}{iq:nw:cross-region-crypto-networks:6}
Race only idempotent messages; for new orders, send once on the best route with a fast fail-over; if a venue remembers keys for the session, race orders
too with one key; never generate a new key for a copy; reconcile every acknowledgement against the key, and cancel any unexpected duplicate fill's
exposure at once.
\iqlookfor{Message classes; venue rule; reconciliation; fail-over rather than racing.}
\end{solution}
